%==============================================================================
\section{Semi-Supervised Learning --- Học máy bán giám sát}
%==============================================================================

\begin{dongco}[title={Giữa supervised và unsupervised}]
Semi-supervised learning không hoàn toàn có giám sát cũng không hoàn toàn không giám sát --- nằm giữa hai phương pháp.
\end{dongco}

\begin{vidu}[title={Minh hoạ semi-supervised}]
\tsimg{11_mo_hinh_va_cac_kieu_hoc/img04.jpg}
\end{vidu}

\begin{ynghia}[title={Bối cảnh sử dụng}]
Huấn luyện trên dataset gồm cả dữ liệu \textbf{có nhãn} và \textbf{không nhãn}.
Thường dùng khi có rất nhiều dữ liệu không nhãn.
Gán nhãn thủ công trong supervised tốn kém; unsupervised thuần có ứng dụng hạn chế $\Rightarrow$ semi-supervised cân bằng hai bên.
\end{ynghia}

\subsection{Semi-Supervised Learning là gì?}

\begin{dinhnghia}[title={Định nghĩa}]
Kỹ thuật ML kết hợp supervised và unsupervised: huấn luyện trên \textbf{ít} dữ liệu có nhãn và \textbf{nhiều} dữ liệu không nhãn.
\end{dinhnghia}

\begin{ynghia}[title={target}]
Phát triển thuật toán chia toàn bộ dữ liệu thành các cụm; điểm gần nhau trong không gian đặc trưng có khả năng cao cùng nhãn đầu ra, rồi gán cụm vào category định sẵn.
\end{ynghia}

\begin{vidu}[title={Tóm tắt}]
\begin{itemize}
  \item Cách tiếp cận ML kết hợp supervised và unsupervised
  \item Huấn luyện bằng dữ liệu có nhãn và không nhãn
  \item Phục vụ classification và regression
\end{itemize}
\end{vidu}

\subsection{So với Supervised Learning}

\begin{ynghia}[title={Khác biệt chính: dataset}]
\textbf{Supervised}: mỗi mẫu có cặp feature--nhãn đích $\Rightarrow$ dự đoán/phân loại chính xác hơn khi nhãn đủ tốt.

\textbf{Semi-supervised}: ít mẫu có nhãn, nhiều mẫu không nhãn; huấn luyện ban đầu trên phần có nhãn, dùng insight đó để học thêm pattern từ phần không nhãn.
\end{ynghia}

\subsection{So với Unsupervised Learning}

\begin{ynghia}[title={Khác biệt}]
\textbf{Unsupervised}: chỉ dữ liệu không nhãn, gom cụm theo feature chung.

\textbf{Semi-supervised}: hỗn hợp có nhãn + không nhãn; hiệu quả hơn vì mỗi cụm có thể được gán nhãn định sẵn nhờ phần dữ liệu đã gán nhãn.
\end{ynghia}

\subsection{Khi nào chọn Semi-Supervised?}

\begin{luuy}[title={Tình huống điển hình}]
Khó và đắt có đủ dữ liệu có nhãn, nhưng thu thập dữ liệu không nhãn dễ hơn.
Khi đó supervised thuần hoặc unsupervised thuần có thể không đủ tốt $\Rightarrow$ dùng semi-supervised.
\end{luuy}

\subsection{Cách Semi-Supervised Learning hoạt động}

\begin{phuongphap}[title={Thành phần huấn luyện}]
Thành phần supervised nhỏ (ít dữ liệu có nhãn) + thành phần unsupervised lớn (nhiều dữ liệu không nhãn).
\end{phuongphap}

\begin{phuongphap}[title={Hai hướng triển khai}]
\begin{enumerate}
  \item Huấn luyện mô hình supervised trên ít dữ liệu có nhãn; áp dụng lên dữ liệu không nhãn để tạo thêm nhãn; huấn luyện lại và lặp.
  \item Dùng unsupervised cluster mẫu tương tự, gán nhãn cho cụm, huấn luyện mô hình với thông tin kết hợp.
\end{enumerate}
\end{phuongphap}

\begin{luuy}[title={Dữ liệu không nhãn phải liên quan}]
Dữ liệu không nhãn phải phù hợp tác vụ mô hình đang học.
Trong ký hiệu xác suất: phân phối đầu vào $p(x)$ phải chứa thông tin về $p(y|x)$ --- xác suất mẫu $x$ thuộc lớp $y$.
\end{luuy}

\subsection{Các giả định (assumptions)}

\subsubsection{Giả định mượt mà (Smoothness Assumption)}

\begin{ynghia}[title={Smoothness Assumption}]
Hai điểm $x_1$, $x_2$ trong vùng mật độ cao (cùng cụm) thì nhãn $y_1$, $y_2$ nên gần nhau.
Ở vùng mật độ thấp, nhãn đầu ra không nhất thiết gần.
\end{ynghia}

\subsubsection{Giả định cụm (Cluster Assumption)}

\begin{ynghia}[title={Cluster Assumption}]
Điểm trong cùng cụm có khả năng cùng lớp.
Dữ liệu không nhãn giúp xác định biên cụm chính xác hơn; dữ liệu có nhãn gán lớp cho từng cụm.
\end{ynghia}

\subsubsection{Giả định mật độ thấp (Low Density Separation)}

\begin{ynghia}[title={Low Density Separation}]
Biên quyết định nên nằm ở vùng mật độ thấp.
Ví dụ nhận dạng chữ số 0 vs 1: điểm ngay trên biên trông như số ``dài'' lẫn lộn --- xác suất gặp mẫu như vậy rất nhỏ.
\end{ynghia}

\subsubsection{Giả định manifold (Manifold Assumption)}

\begin{ynghia}[title={Manifold Assumption}]
Trong không gian đầu vào chiều cao, dữ liệu nằm trên các manifold chiều thấp hơn; điểm cùng nhãn nằm trên cùng manifold.
\end{ynghia}

\subsection{Kỹ thuật Semi-Supervised}

\begin{phuongphap}[title={Một số kỹ thuật phổ biến}]
\begin{itemize}
  \item \textbf{Self-training}: chỉnh supervised (classification/regression) để dùng cả nhãn và không nhãn.
  \item \textbf{Co-training}: mở rộng self-training với nhiều ``view'' dữ liệu (ví dụ trang web: nội dung text + hyperlink); hai classifier riêng trên hai view cải thiện học.
  \item \textbf{Graph-based label propagation}: mô hình hoá dữ liệu thành đồ thị (node = điểm, cạnh = độ tương đồng); nhãn lan truyền từ điểm có nhãn sang láng giềng; cập nhật lặp để gán nhãn cho node chưa nhãn.
\end{itemize}
\end{phuongphap}

\subsection{Thách thức}

\begin{luuy}[title={Thách thức}]
\begin{itemize}
  \item \textbf{Chất lượng dữ liệu không nhãn}: nhiễu hoặc không liên quan $\Rightarrow$ dự đoán sai.
  \item \textbf{Distribution shift}: dữ liệu có nhãn và không nhãn khác phân phối (ví dụ ảnh rõ vs ảnh camera an ninh) $\Rightarrow$ khó tổng quát hoá.
\end{itemize}
\end{luuy}

\subsection{Ứng dụng}

\begin{phuongphap}[title={Lĩnh vực}]
Text classification, image classification, speech analysis, anomaly detection --- target chung: gán thực thể vào category định sẵn; điểm gần nhau có khả năng cùng nhãn.

\textbf{Speech Recognition}: gán nhãn audio tốn thời gian; kết hợp audio không nhãn với ít bản ghi có transcript.

\textbf{Web Content Classification}: hàng tỷ trang, không thể gán nhãn thủ công; hỗ trợ công cụ tìm kiếm xếp hạng nội dung.

\textbf{Text Document Classification}: học từ ít tài liệu có nhãn + corpus lớn không nhãn $\Rightarrow$ tăng độ chính xác không cần dataset nhãn khổng lồ.
\end{phuongphap}
